\documentclass[10pt]{article}
\usepackage[utf8]{inputenc}
\usepackage[T1]{fontenc}
\usepackage{amsmath}
\usepackage{amsfonts}
\usepackage{amssymb}
\usepackage[version=4]{mhchem}
\usepackage{stmaryrd}
\usepackage{bbold}
\usepackage{graphicx}
\usepackage[export]{adjustbox}
\graphicspath{ {./images/} }

\begin{document}
Mechanism design

\section*{Some references}
See Myerson (1981), Mas-Colell et al. (1995), Krishna (2010), and Menezes and Monteiro (2005).

\section*{Mechanism design in IPVM}
\begin{itemize}
  \item Goal. What is the best way to transfer the object?
  \item best? Objective function. Among possibilities:
  \item seller's expected revenue.
  \item aggregate surplus (buyers surplus).
  \item way? Mechanisms. Problem: too many.
  \item How to formalize the notion of mechanism so as to cover as many as possible transaction forms while still keeping the optimization problem tractable?
  \item How to solve the optimization program?
\end{itemize}

\section*{Definitions}
Definition 1 (Environment) An environment is an $E =\left[O,\left\{\left(X_{i}, X_{i}, \mu_{i}\right), V_{i}\right\}_{i=1}^{I}\right]$ where $O$ is a set of outcomes, $\left(X_{i}, X_{i}, \mu_{i}\right)$ is the space of agent $i$ 's private information, and $V_{i}: X \times O \rightarrow \mathbb{R}$ is i's payoffs given a realization of information and an outcome.

Definition 2 (Mechanism) Given $E$, a mechanism (or game form) is an $M =\left\{\left\{S_{i}\right\}_{i=1}^{I}, g: \prod_{i=1}^{I} S_{i} \rightarrow O\right\}$ where $S_{i}$ is player $i$ 's set of available messages or signals, and $g$ is a function, termed the outcome function, that assigns to each message profile an outcome.

\section*{Environment}
\begin{center}
\includegraphics[max width=\textwidth, alt={}]{9c2e101e-e880-4708-a228-aa1599e90ecf-06_744_2569_351_153}
\end{center}

\section*{Environment and mechanism}
\begin{center}
\includegraphics[max width=\textwidth, alt={}]{9c2e101e-e880-4708-a228-aa1599e90ecf-07_1467_2573_355_153}
\end{center}

\begin{itemize}
  \item Given an environment $E$, every mechanism $M$ induces a Bayesian game $\Gamma_{M}=\left\{\left(X_{i}, X_{i}, \mu_{i}\right), S_{i}, u_{i}\right\}_{i=1}^{I}$,
  \item Player $i$ observes private information $x_{i} \in X_{i}$ and chooses an action $s_{i} \in S_{i}$.\\
A strategy for player $i$ is a function $\phi_{i}: X_{i} \rightarrow S_{i}$
  \item Given an action profile $s, g(s) \in O$ is the outcome of the game.
  \item Player $i$ 's payoff is $u_{i}(x, s)=V_{i}(x, g(s))$.
\end{itemize}

$$
u_{i}\left(x_{1}, \ldots, x_{n}, q_{1}, q_{2}, \ldots, q_{n}, m_{1}, \ldots, m_{n}\right)=x_{i} \cdot q_{i}-m_{i}
$$

\section*{Direct mechanism}
Definition 3 (Direct mechanism) A mechanism $M=\left[\left\{S_{i}\right\}_{i=1}^{I}, g\right]$ is direct if $S_{i}=X_{i}$ for every $i$.

\begin{itemize}
  \item Bidders observe their types and report $x^{\prime}$ individually and independently.
  \item Given a report profile $x^{\prime}=\left(x_{1}^{\prime}, \ldots, x_{I}^{\prime}\right)$, the function $g$ determines the outcome $g\left(x^{\prime}\right) \in O$.
\end{itemize}

Definition 4 (Direct revelation mechanism) If the game $\Gamma_{M}$ induced by a direct mechanism $M$ has an equilibrium $\left\{\phi_{i}\right\}_{i=1}^{I}$ where for every $i$ and every $x_{i} \in X_{i}, \phi\left(x_{i}\right)=x_{i}$, then $M$ is a direct revelation mechanism (DRM).

\begin{itemize}
  \item BNE versus DSE.
\end{itemize}

\section*{Example: application to the IPVM}
\begin{itemize}
  \item $O=\left\{(q, m): q \in[0,1]^{I}, \sum_{i=1}^{I} q_{i} \leq 1, m \in \mathbb{R}^{I}\right\}$
  \item A mechanism $M=\left[\left\{S_{i}\right\}_{i=1}^{I}, g: \prod_{i=1}^{I} S_{i} \rightarrow O\right]$ :
  \item A signaling space $S_{i}$ per agent, $s_{i} \in S_{i}, s=\left(s_{1}, s_{2}, \ldots, s_{I}\right)$.
  \item $g(s)=(q(s), t(s))$,
  \item $q(s)=\left(q_{1}(s), \ldots, q_{I}(s)\right), \sum_{i=1}^{I} q(s) \leq 1$,
  \item $t(s)=\left(t_{1}(s), \ldots, t_{I}(s)\right)$.
\end{itemize}

Example FPA: mechanisms that is not direct, $x_{i} \in[0,1]$

\begin{itemize}
  \item $S_{i}=\mathbb{R}, b_{i} \in \mathbb{R}$ is $i$ 's bid.
  \item $g(b)=g\left(b_{1}, \ldots, b_{I}\right)=\left\{q_{i}(b), t_{i}(b)\right\}_{i=1}^{I}$
  \item $q_{i}(b)= \begin{cases}1 & \text { if } b_{i}>b_{-i}^{(1)} \\ 0 & \text { if } b_{i}<b_{-i}^{(1)} \\ . & \text { if } b_{i}=b_{-i}^{(1)}\end{cases}$
  \item $t_{i}(b)=b_{i} q_{i}(b)$.
\end{itemize}

Example SPA: (it can be seen as a direct mechanism)

\begin{itemize}
  \item $S_{i}=X_{i}=[0,1], b_{i} \in X_{i}=[0,1]$ is $i$ 's report (and bid).
  \item $q_{i}(b)$ as in the FPA.
  \item $t_{i}(b)= \begin{cases}b_{-i}^{(1)} & \text { if } b_{-i}^{(1)}<b_{i} \\ 0 & \text { if } b_{-i}^{(1)}>b_{i} \\ \cdot & \text { if } b_{-i}^{(1)}=b_{i}\end{cases}$
\end{itemize}

\section*{Revelation principle}
Theorem 1 (Revelation principle) Let $\phi$ be an equilibrium of $\Gamma_{M}$, the game induced by a mechanism $M$. Then there exists an equilibrium $\phi^{\prime}$ of (the game induced by) a direct mechanism $g^{\prime}$ such that the $\phi^{\prime}$ and $\phi$ generate the same outcome and for every $\forall i \wedge x_{i} \in X_{i}, \phi_{i}^{\prime}\left(x_{i}\right)=x_{i}$.

\section*{Proof 1}
\begin{center}
\includegraphics[max width=\textwidth, alt={}]{9c2e101e-e880-4708-a228-aa1599e90ecf-14_655_2380_387_169}
\end{center}

\section*{Proof 2}
\begin{center}
\includegraphics[max width=\textwidth, alt={}]{9c2e101e-e880-4708-a228-aa1599e90ecf-15_1233_2380_387_169}
\end{center}

\section*{Proof 3}
\begin{center}
\includegraphics[max width=\textwidth, alt={}]{9c2e101e-e880-4708-a228-aa1599e90ecf-16_1233_2380_387_169}
\end{center}

\section*{Proof 4}
\begin{center}
\includegraphics[max width=\textwidth, alt={}]{9c2e101e-e880-4708-a228-aa1599e90ecf-17_1237_2400_383_149}
\end{center}

\section*{Proof 5}
\begin{center}
\includegraphics[max width=\textwidth, alt={}]{9c2e101e-e880-4708-a228-aa1599e90ecf-18_1245_2396_367_153}
\end{center}

Proof(Sketch).

\begin{itemize}
  \item Define DRM
  \item $(\forall i) S_{i}=X_{i}=[0,1]$.
  \item $\left(\forall i, \forall x \in X=[0,1]^{I}\right)$
\end{itemize}

$$
q_{i}^{\prime}\left(x_{-i}, x_{i}\right)=q_{i}\left(\phi_{1}\left(x_{1}\right), \cdots, \phi_{I}\left(x_{I}\right)\right)
$$

and

$$
t_{i}^{\prime}\left(x_{-i}, x_{i}\right)=t_{i}\left(\phi_{-i}\left(x_{-i}\right), \phi_{i}\left(x_{i}\right)\right)
$$

\begin{itemize}
  \item Check that $\left(q^{\prime}, t^{\prime}\right)$ generate the same outcome as $(q, t)$ and that truthtelling is a BNE under ( $q^{\prime}, t^{\prime}$ ) to complete the proof.
\end{itemize}

\section*{Modeling question}
\begin{itemize}
  \item Modeling question: model of mechanisms must be broad and tractable.
  \item Revelation Principle implies two simplifications
\end{itemize}

\begin{enumerate}
  \item We may restrict attention to direct mechanisms, message space is space of private information.
  \item We may farther restrict attention to direct mechanisms that have truthful revelation in equilibrium.
\end{enumerate}

\begin{itemize}
  \item It is clear how to use (1); how does one take advantage of (2)?
  \item Remark. Careful with use of Revelation Principle.
\end{itemize}

\section*{References}
Krishna, Vijay. 2010. Auction Theory. 2nd ed. Elsevier, Academic Press.\\
Mas-Colell, Andreu, Michael Whinston, and Jerry Green. 1995. Microeconomic Theory. Oxford University Press.\\
Menezes, Flavio M., and Paulo K. Monteiro. 2005. An Introduction to Auction Theory. Oxford University Press.\\
Myerson, Roger. 1981. "Optimal Auction Design." Mathematics of Operations Research 6: 58-73.


\end{document}